\documentclass[10pt,a4paper]{amsart}
\usepackage[margin=19mm]{geometry}
\usepackage{amsmath,amssymb,mathtools}
\usepackage[scr=rsfs,cal=euler]{mathalfa}
\usepackage{xfrac}
\usepackage[unicode,hidelinks,bookmarks=false]{hyperref}
\newcommand{\ie}{i.e.\ }
\newcommand{\cf}{cf.\ }
\newcommand{\resp}{respectively\ }
\newcommand{\Subsection}{\subsection}
\providecommand{\nobreakdash}{\nobreak\hbox{-}}
\setlength{\emergencystretch}{2em}
\multlinegap0pt
\usepackage{bm}
\usepackage{bbm}
\usepackage{dsfont}
\usepackage{xspace}
\usepackage{cancel}
\usepackage{mathtools}
\usepackage{amsthm}
\makeatletter
\@ifundefined{theorem}{\newtheorem{theorem}{Theorem}}{}
\@ifundefined{lemma}{\newtheorem{lemma}[theorem]{Lemma}}{}
\@ifundefined{proposition}{\newtheorem{proposition}[theorem]{Proposition}}{}
\makeatother
\title[Strictly stable solutions in uniformly convex planar domains]{Strictly stable solutions in uniformly convex planar domains may have nonconvex superlevel sets}
\author{Yi Ru-Ya Zhang}
\date{Source: arXiv:2607.06031v1 (CC BY 4.0)}
\begin{document}
\maketitle

\noindent\emph{Source and licence.} Strictly stable solutions in uniformly convex planar domains may have nonconvex superlevel sets. Version arXiv:2607.06031v1 (CC BY 4.0), distributed under the Creative Commons Attribution 4.0 International licence, as recorded in the arXiv licence metadata for this exact version and shown on the abstract page https://arxiv.org/abs/2607.06031v1. Full rights notice: licenses/arxiv-2607.06031-CC-BY-4.0.txt.\par\smallskip

\noindent\textit{Editorial scope.} Counted unit: the Proposition whose source label is \texttt{prop:c0-approx} in the article (source0.tex lines 1576--1593, source label \texttt{prop:c0-approx}), delivered together with the article's own complete original proof of it (source0.tex lines 1595--1747, one \texttt{proof} environment of 153 lines). No mathematics is written, repaired or re-derived: statement and proof are the article's own text line for line. Required context retained uncounted: eq:defn-mu-star (lines 1230--1232); lem:slow-channel-maximum (lines 1438--1456); lem:slow-profile-xbounds (lines 1496--1519). Every definition, display and label of those lines is retained unchanged. Omitted from this package and not redistributed: 1--1229, 1233--1437, 1457--1495, 1520--1575, 1594--1594, 1748--2285. Nothing inside a retained excerpt is omitted or altered: every retained body is delivered exactly as the article's lines are written, so no omission receipt of any kind accompanies this package. Citations: the retained excerpts carry no citation command at all, so no bibliography entry of the article is transcribed and no reference is redistributed. Reference adaptation: none was needed and none was made; every reference inside the delivered unit resolves inside it, so no unresolved reference or citation is produced. Printed slips kept literal: no printed word, symbol, hypothesis or number of the counted statement or of its proof was altered. Layout and class: the article's own class is \texttt{amsart} with packages including a4wide, color, eucal, enumerate, mathrsfs, ulem, amsmath, amssymb, epsfig, amsthm, inputenc, psfrag, hyperref, cleveref; the wrapper is the lane's pinned \texttt{amsart} editorial wrapper at 10pt on A4, which restates the theorem-like environments the retained text needs and adds the article's own macro definitions that the retained text needs (none), with extra wrapper packages (bm, bbm, dsfont, xspace, cancel, mathtools). The delivered numbering is the unit's own. Assets: no retained span contains a figure environment, a graphics-inclusion command, a PDF-page inclusion, a table or a TikZ picture; no image file is copied into this package, the package declares no asset dependency and no image file is redistributed with it. Licence: version arXiv:2607.06031v1 of this article is distributed under the Creative Commons Attribution 4.0 International licence, as recorded in the arXiv licence metadata for this exact version and shown on the abstract page https://arxiv.org/abs/2607.06031v1. A full rights notice is delivered at licenses/arxiv-2607.06031-CC-BY-4.0.txt. Characters: every retained excerpt is pure ASCII, so the delivered engine, XeLaTeX with its default Latin Modern text font, reports no missing character.
\begin{equation}\label{eq:defn-mu-star}
    \nu_*:=\lambda_1\left(-\frac{d^2}{d\xi^2}-f'(\overline U),\,(-h_*,h_*)\right)
\end{equation}

\begin{lemma} \label{lem:slow-channel-maximum}
Let $c\in L^\infty(\mathcal Q_I)$ satisfy
$$
0\le c(X,\eta)\le f'(\overline U(\eta))
\qquad\text{for a.e. }(X,\eta)\in\mathcal Q_I.
$$
Let $w\in W^{1,2}(\mathcal Q_I)\cap C(\overline{\mathcal Q_I})$ satisfy
$$
-\epsilon^2 w_{XX}-w_{\eta\eta}-c(X,\eta)w\ge0
$$
in the weak sense in $\mathcal Q_I$, and assume that
$$
w\ge0\qquad\text{on }\partial\mathcal Q_I.
$$
Then
$$
w\ge0\qquad\text{in }\mathcal Q_I.
$$
\end{lemma}

\begin{lemma}
\label{lem:slow-profile-xbounds}
Let $K=[h_1,h_2]\subset\subset(0,h_c)$ be a compact interval of half-widths contained in the stable lower branch. For $h\in K$, let $\mathcal U_h$ denote the lower-branch profile on $(-h,h)$, and define
$$
W(h,\sigma):=\mathcal U_h(h\sigma),
\qquad
(h,\sigma)\in K\times[-1,1].
$$
Then $W\in C^2(K\times[-1,1])$. 

Consequently, if $I\subset\subset J$, 
$H\in C^2(\overline I)$,   $H(\overline I)\subset K$, and 
$$
V(X,\eta):=\mathcal U_{H(X)}(\eta),
\qquad
(X,\eta)\in\mathcal Q_I:=\{(X,\eta):X\in I,\ |\eta|<H(X)\},
$$
then there exists a constant $C_V>0$, depending only on $f$, $K$, and
$\|H\|_{C^2(\overline I)}$, such that
$$
|V_{XX}(X,\eta)|\le C_V
\qquad\text{for every }(X,\eta)\in\mathcal Q_I.
$$
\end{lemma}

\begin{proposition}\label{prop:c0-approx}
Let $I'\subset\subset I$.  There exist constants
$$
C>0,
\qquad
\epsilon_0>0,
\qquad
\kappa>0,
$$
depending only on $f$, on $h_*$, on the compact stable branch interval
containing $H(\overline I)$, on $\|H\|_{C^2(\overline I)}$, and on the
intervals $I'\subset\subset I$, such that, for every $\epsilon\in(0,\epsilon_0]$,
$$
\sup_{(X,\eta)\in\mathcal Q_{I'}}
|v_\epsilon(X,\eta)-V(X,\eta)|
\le C\epsilon^2.
$$
\end{proposition}

\begin{proof}
Set $W_\epsilon:=v_\epsilon-V$
on $\mathcal Q_I$.  Since $V$ solves the one-dimensional equation in the
vertical variable,
$$
-\epsilon^2V_{XX}-V_{\eta\eta}-f(V)=-\epsilon^2V_{XX}.
$$
Subtracting the equation for $V$ from the equation for $v_\epsilon$, we
obtain
\begin{equation}\label{eq:difference-equation}
-\epsilon^2 (W_\epsilon)_{XX}-(W_\epsilon)_{\eta\eta}-c_\epsilon(X,\eta)W_\epsilon
=\epsilon^2 V_{XX}
\qquad\text{in }\mathcal Q_I,
\end{equation}
where
$$
c_\epsilon(X,\eta):=
\begin{cases}
\dfrac{f(v_\epsilon(X,\eta))-f(V(X,\eta))}{v_\epsilon(X,\eta)-V(X,\eta)},
& v_\epsilon(X,\eta)\ne V(X,\eta),\\
f'(V(X,\eta)),& v_\epsilon(X,\eta)=V(X,\eta).
\end{cases}
$$
Since $0\le v_\epsilon,V\le\overline U$ and $f'$ is nondecreasing,
\begin{equation}\label{eq:ceps-bound}
    0\le c_\epsilon(X,\eta)\le f'(\overline U(\eta))
\qquad\text{in }\mathcal Q_I.
\end{equation}

 

% Since $H(\overline I)$ is contained in a compact subinterval of the stable lower branch and $H$ is smooth, we claim that the map $$(X,\eta)\mapsto U_{H(X)}(\eta)$$
% has bounded second derivative in $X$ on $\overline{\mathcal Q_I}$.  
% Indeed, write $\eta=h\sigma$, with $|\sigma|\le1$, and consider
% $(h,\sigma)\mapsto U_h(h\sigma)$ on the compact set of stable widths
% $h\in H(\overline I)$. Note that, the smooth dependence of ordinary differential equations
% and invertibility of the lower-branch map give uniform bounds for the required
% $h$-derivatives, and then the chain rule for $h=H(X)$  gives the claimed $X$-bounds.  Hence
% there is a constant $C_V>0$, depending only on $f$, the chosen stable branch
% interval, and $H|_{\overline I}$, such that
% $$
% |V_{XX}(X,\eta)|\le C_V
% \qquad\text{for every }(X,\eta)\in\mathcal Q_I.
% $$

Choose a compact  interval
$K\subset\subset(0,h_c)$ such that $ H(\overline I)\subset K.$
By Lemma~\ref{lem:slow-profile-xbounds}, there exists a constant $C_V>0$,
depending only on $f$, $K$, and $\|H\|_{C^2(\overline I)}$, such that
\begin{equation}\label{eq:VXX-bound}
|V_{XX}(X,\eta)|\le C_V
\qquad\text{for every }(X,\eta)\in\mathcal Q_I.
\end{equation}

Now we start to construct barrier functions for $W_\epsilon$. 
Let $\varphi_*$ be the positive first Dirichlet eigenfunction of the operator
$$-\frac{d^2}{d\xi^2}-f'(\overline U)$$
on $(-h_*,h_*)$, normalized by $\max_{[-h_*,h_*]}\varphi_*=1. $
Thus, recalling \eqref{eq:defn-mu-star},
$$
-\varphi_*''-f'(\overline U)\varphi_*=\nu_*\varphi_*,
\qquad
\varphi_*>0\text{ in }(-h_*,h_*),
\qquad
\varphi_*(\pm h_*)=0.
$$
Since $\sup_{X\in\overline I}H(X)<h_*,$ then
there is a constant $m_\varphi>0$, depending only on $H|_{\overline I}$ and $h_*$, such that
$$
\varphi_*(\eta)\ge m_\varphi
\qquad\text{for every }(X,\eta)\in\mathcal Q_I.
$$
Using $c_\epsilon\le f'(\overline U)$, we have
$$
-\varphi_*''-c_\epsilon\varphi_*
= \nu_*\varphi_*+\bigl(f'(\overline U)-c_\epsilon\bigr)\varphi_*
\ge \nu_*\varphi_*.
$$
Choose $\kappa>0$ such that $0<\kappa^2<\nu_*.$
Define
$$
E_L(X):=e^{-\kappa(X-X_L)/\epsilon},
\qquad
E_R(X):=e^{-\kappa(X_R-X)/\epsilon}.
$$
Then
$$
\left(-\epsilon^2\partial_{XX}-\partial_{\eta\eta}-c_\epsilon\right)(E_L\varphi_*)
=E_L\left(-\kappa^2\varphi_*-\varphi_*''-c_\epsilon\varphi_*\right)
\ge 0,
$$
and the same inequality holds for $E_R\varphi_*$.

Let
\begin{equation}\label{eq:defn-M-star}
    M_*:=\|\overline U\|_{L^\infty((-h_*,h_*))}.
\end{equation}
Choose constants $C_0$ and $C_1$, depending only on the data listed in the
statement, so large that
$$
C_0\nu_* m_\varphi\ge C_V,
\qquad
C_1 m_\varphi\ge 2M_*.
$$
Set
$$
\mathcal B(X,\eta):=
C_0\epsilon^2\varphi_*(\eta)
+C_1\bigl(E_L(X)+E_R(X)\bigr)\varphi_*(\eta).
$$
The inequalities above imply
$$
\left(-\epsilon^2\partial_{XX}-\partial_{\eta\eta}-c_\epsilon\right) \mathcal B
\ge \epsilon^2 C_V
\ge \epsilon^2 \left| V_{XX}\right|
\qquad\text{in }\mathcal Q_I.
$$
Note that, on the top and bottom boundary arcs $\eta=\pm H(X)$, both $v_\epsilon$ and
$V$ vanish, and thus $W_\epsilon=0$.  On the artificial vertical sides $X=X_L$ and $X=X_R$,
the estimate $|W_\epsilon|\le2 M_*$ holds by \eqref{eq:defn-M-star}, while $\mathcal B\ge C_1m_\varphi\ge2 M_*$.
Therefore
$$
-\mathcal B\le W_\epsilon\le\mathcal B
\qquad\text{on }\partial\mathcal Q_I.
$$

Now the comparison principle in Lemma~\ref{lem:slow-channel-maximum} applies
due to \eqref{eq:ceps-bound}.
Applying that lemma to $\mathcal B-W_\epsilon$ and to
$\mathcal B+W_\epsilon$, and using \eqref{eq:difference-equation}, gives
$$
|W_\epsilon(X,\eta)|\le\mathcal B(X,\eta)
\qquad\text{in }\mathcal Q_I.
$$

Let $d_I:=\operatorname{dist}(I',\partial I)>0.$
For $X\in I'$,
$$
E_L(X)+E_R(X)
\le 2e^{-\kappa d_I/\epsilon}.
$$
After decreasing $\epsilon_0>0$ if necessary, we have
$$
e^{-\kappa d_I/\epsilon}\le \epsilon^2
\qquad\text{for }\ 0<\epsilon\le\epsilon_0.
$$
Consequently,
$$
|W_\epsilon(X,\eta)|\le C\epsilon^2
\qquad\text{for every }(X,\eta)\in\mathcal Q_{I'},
$$
where $C$ depends only on the stated data.  This proves the proposition.
\end{proof}

\end{document}
